\documentclass[12pt,a4paper]{article}
\usepackage[UTF8]{ctex}
\usepackage{amsmath,amssymb,bm}
\usepackage{geometry}
\usepackage{hyperref}
\usepackage{booktabs}
\geometry{margin=1in}

\title{simulateTimeVaryingFirDouble 与 simulateTimeVaryingFirFixed 说明文档}
\author{}
\date{\today}

\begin{document}
\maketitle

\section{文档目的}

本文档用于详细说明 \texttt{FirConvValidationTool.m} 中以下两个函数的实现思路与处理流程：

\begin{itemize}
\item \texttt{simulateTimeVaryingFirDouble}
\item \texttt{simulateTimeVaryingFirFixed}
\end{itemize}

这两个函数分别对应：

\begin{itemize}
\item 浮点参考链路的时变 FIR / MIMO 卷积
\item 定点链路的时变 FIR / MIMO 卷积
\end{itemize}

两者共享同一个总体目标：在给定输入 IQ 数据与时变信道系数的情况下，按时间顺序计算每个输出通道的复数卷积结果。

\section{问题背景}

\subsection{为什么不是普通的固定系数 FIR}

普通 FIR 可以写成固定系数卷积：

\begin{equation}
y[n] = \sum_{k=0}^{K-1} h[k]\,x[n-k]
\end{equation}

其中系数 $h[k]$ 在整个计算过程中不变化。

但本项目中的信道是\textbf{时变}的。信道矩阵 $H$ 由多个 CIR snapshot 组成，每个 snapshot 都包含一组 tap：

\begin{equation}
H(s, t, ch)
\end{equation}

其中：

\begin{itemize}
\item $s$ 表示 snapshot 索引
\item $t$ 表示 tap 索引
\item $ch$ 表示输入输出通道组合后的信道索引
\end{itemize}

因此这里不能简单地把整个问题看成一次固定系数卷积，而必须在输出时间轴上逐点判断当前属于哪个 snapshot，然后使用该 snapshot 对应的 tap 系数和 delay。

\subsection{MIMO 结构}

若输入通道数为 $IN\_num$，输出通道数为 $OUT\_num$，则一共有

\begin{equation}
IN\_num \times OUT\_num
\end{equation}

个子信道 FIR 单元。

项目中采用的通道索引映射关系为：

\begin{equation}
ch = (m - 1)\cdot OUT\_num + n
\end{equation}

其中：

\begin{itemize}
\item $m$ 是输入通道编号
\item $n$ 是输出通道编号
\end{itemize}

对每个输出通道，需要将所有输入通道对应的 FIR 输出结果求和。

\section{统一数学模型}

\subsection{时变复数 FIR 的输出公式}

设：

\begin{itemize}
\item $x_m[n]$ 为第 $m$ 路输入 IQ
\item $y_n[n]$ 为第 $n$ 路输出 IQ
\item $h_{m,n}^{(s)}[t]$ 为 snapshot $s$ 时，第 $(m,n)$ 个子信道第 $t$ 个 tap 的复数系数
\item $d_{m,n}^{(s)}[t]$ 为对应的 delay
\end{itemize}

则输出可写为：

\begin{equation}
y_n[k]
=
\sum_{m=1}^{IN\_num}
\sum_{t=1}^{T\_num}
h_{m,n}^{(s(k))}[t]\,
x_m\!\left[k - d_{m,n}^{(s(k))}[t]\right]
\end{equation}

其中：

\begin{equation}
s(k) = \min\left(\left\lfloor \frac{k-1}{N_{\text{snap}}} \right\rfloor + 1,\; N_{\text{samples}}\right)
\end{equation}

$N_{\text{snap}}$ 表示每个 snapshot 对应的采样点数：

\begin{equation}
N_{\text{snap}} = \mathrm{round}\!\left(\frac{f_s}{CIR\_update\_rate}\right)
\end{equation}

这正是代码中 \texttt{samplesPerSnapshot} 的来源。

\section{simulateTimeVaryingFirDouble 说明}

\subsection{作用}

\texttt{simulateTimeVaryingFirDouble} 用于生成\textbf{浮点参考输出}。

它的核心思想是：

\begin{itemize}
\item 输入 IQ 使用 \texttt{double complex}
\item 信道系数也使用浮点复数
\item delay 已按当前阶段规则换算成样值移位
\item 不考虑中途量化误差
\end{itemize}

因此它代表的是理想参考链路。

\subsection{输入含义}

函数输入：

\begin{itemize}
\item \texttt{inputs}：尺寸为 $[N, IN\_num]$ 的复数输入矩阵
\item \texttt{H.coeff}：尺寸为 $[N_{snap}, T\_num, IN\_num\cdot OUT\_num]$ 的复数系数
\item \texttt{H.delay\_samples}：与系数同维度的样值延迟
\item \texttt{samplesPerSnapshot}：每个 snapshot 持续的采样点数
\end{itemize}

\subsection{流程拆解}

\subsubsection{步骤 1：确定输出长度}

输出长度设置为：

\begin{equation}
N_{\text{out}} = N_{\text{in}} + d_{\max}
\end{equation}

其中 $d_{\max}$ 是所有 tap 中最大的 delay。

这样做的原因是：信号经过延迟后，尾部会比输入更长，因此需要预留额外长度容纳最后一段能量。

\subsubsection{步骤 2：按输出时间索引逐点计算}

代码主循环的外层变量是输出样点索引 $n$：

\begin{equation}
n = 1,2,\dots,N_{\text{out}}
\end{equation}

对每个 $n$：

\begin{enumerate}
\item 先判断当前属于哪个 snapshot
\item 再对所有输入通道和输出通道组合进行求和
\item 再对该子信道下所有 tap 求和
\end{enumerate}

\subsubsection{步骤 3：根据 snapshot 选择系数}

在输出样点 $n$ 上，snapshot 通过下面的规则确定：

\begin{equation}
s = \min\left(\left\lfloor \frac{n-1}{samplesPerSnapshot}\right\rfloor + 1,\; N_{snap}\right)
\end{equation}

这对应“零阶保持”逻辑：在一个 snapshot 时间窗内，系数保持不变。

\subsubsection{步骤 4：按 delay 对输入取样}

对某个 tap，如果当前 snapshot 下的 delay 为 $d$，则输入索引为：

\begin{equation}
srcIdx = n - d
\end{equation}

若 $srcIdx$ 越界，则该 tap 在此时刻贡献为 0，这等价于 FIR 的零填充边界处理。

\subsubsection{步骤 5：复数乘法并累加}

若输入样点为 $x$，系数为 $h$，则贡献为：

\begin{equation}
\Delta y = x \cdot h
\end{equation}

所有 tap 的贡献先在单个子信道内求和，再跨输入通道求和，最终得到每个输出通道的输出。

\subsection{为什么采用“按时间点直接求和”}

虽然文档中提到可使用 MATLAB 自带卷积函数，但在当前实现中采用了更直接的按时间索引求和方式。原因是：

\begin{itemize}
\item 当前信道是时变的，不是固定系数
\item 每个 snapshot 的 delay 和系数都可能变化
\item 直接按时间索引求值更容易保证逻辑正确
\item 更接近 FPGA 真实按时钟推进的思维方式
\end{itemize}

\section{simulateTimeVaryingFirFixed 说明}

\subsection{作用}

\texttt{simulateTimeVaryingFirFixed} 用于生成\textbf{定点链路下的高精度累加结果}。

注意，这个函数的输出并不是最终 16 bit IQ，而是：

\begin{itemize}
\item 已经完成所有 tap 卷积和多输入通道求和
\item 但还没有执行最终动态截位
\item 保持在较高整数精度下的中间结果
\end{itemize}

这个中间结果随后会交给 shift 搜索模块处理。

\subsection{输入含义}

函数输入：

\begin{itemize}
\item \texttt{inputs}：由 \texttt{.txt} 解析得到的复数 \texttt{int16} 输入
\item \texttt{H.qReal}, \texttt{H.qImag}：由 \texttt{.irc} 恢复出的系数实部/虚部
\item \texttt{H.delay\_clks}：由 \texttt{.ird} 恢复出的 delay 时钟计数
\item \texttt{samplesPerSnapshot}：每个 snapshot 对应的样点数
\end{itemize}

\subsection{为什么不能只保留 32 bit}

项目需求中提到：

\begin{quote}
定点乘法与加法全程保持 32 bit 全精度
\end{quote}

严格来说，这句话只对\textbf{单次乘法}成立。

若输入和系数都是 16 bit，则单次乘法量级约为：

\begin{equation}
16 \text{ bit} \times 16 \text{ bit} \Rightarrow 32 \text{ bit}
\end{equation}

但卷积中会有两类额外增长：

\begin{enumerate}
\item 同一个子信道中多个 tap 的累加
\item 同一个输出通道中多个输入通道结果的累加
\end{enumerate}

例如若有 $T\_num$ 个 tap，则仅 tap 求和就可能增加大约

\begin{equation}
\lceil \log_2(T\_num) \rceil
\end{equation}

位的动态范围。

若再叠加 $IN\_num$ 路输入通道求和，还会进一步增长。

因此代码里统一采用 \texttt{int64} 作为累加容器，这样可以避免在中途发生额外截断。

\subsection{复数乘法展开}

若输入 IQ 为

\begin{equation}
x = x_r + jx_i
\end{equation}

系数为

\begin{equation}
h = c_r + jc_i
\end{equation}

则两者乘积为：

\begin{equation}
xh = (x_r c_r - x_i c_i) + j(x_r c_i + x_i c_r)
\end{equation}

代码中正是按这个公式分别累加实部和虚部：

\begin{align}
accR &\leftarrow accR + (x_r c_r - x_i c_i) \\
accI &\leftarrow accI + (x_r c_i + x_i c_r)
\end{align}

\subsection{流程拆解}

\subsubsection{步骤 1：与浮点链路保持同样的时间推进}

外层仍按输出样点索引 $n$ 推进，并根据 $n$ 计算当前 snapshot。

这保证定点链路和浮点链路在时间轴上严格对齐。

\subsubsection{步骤 2：按 delay\_clk 取输入样点}

当前阶段由于 $clk = f_s$，所以：

\begin{equation}
delay\_clk = delay\_samples
\end{equation}

因此定点链路中可直接用 \texttt{delay\_clks} 作为样值移位。

\subsubsection{步骤 3：单 tap 乘法}

对每个 tap：

\begin{itemize}
\item 读取输入 I/Q 的 \texttt{int16}
\item 读取系数实部/虚部的 \texttt{int16}
\item 用 \texttt{int64} 执行复数乘法展开
\end{itemize}

\subsubsection{步骤 4：tap 内累加}

在同一个子信道内，对所有 tap 的乘法结果累加，得到该子信道在当前输出时刻的贡献。

\subsubsection{步骤 5：跨输入通道累加}

对同一个输出通道，将所有输入通道对应子信道的结果再求和。

最终得到：

\begin{itemize}
\item \texttt{accumReal(n, outIdx)}
\item \texttt{accumImag(n, outIdx)}
\end{itemize}

\subsection{函数输出的物理意义}

\texttt{simulateTimeVaryingFirFixed} 的输出可以理解为：

\begin{quote}
如果 FPGA 在最终输出位宽限制之前，先把所有乘法和加法都保存在一个足够宽的寄存器中，那么这个寄存器里应该接近什么数值。
\end{quote}

所以它是动态截位模块的输入，而不是最终可直接和 FPGA 16 bit 输出逐样比较的值。

\section{两者之间的关系}

\subsection{一个是参考，一个是待量化结果}

\texttt{simulateTimeVaryingFirDouble} 给出的是浮点参考结果：

\begin{equation}
y_{\text{float}}[n]
\end{equation}

\texttt{simulateTimeVaryingFirFixed} 给出的是高精度定点累加结果：

\begin{equation}
y_{\text{acc}}[n]
\end{equation}

然后动态 shift 搜索模块会对 $y_{\text{acc}}[n]$ 施加右移、饱和和反标定，得到用于误差比较的结果：

\begin{equation}
y_{\text{fixed}}^{(N)}[n]
\end{equation}

其中 $N$ 为候选 shift。

\subsection{最终目标}

目标不是让高精度累加值直接接近浮点值，而是要找一个合适的 shift，使得最终 16 bit 输出在误差意义下最接近浮点参考：

\begin{equation}
N^\star = \arg\min_N \; \mathrm{MSE}\left(y_{\text{float}}, y_{\text{fixed}}^{(N)}\right)
\end{equation}

\section{更直观的理解方式}

\subsection{浮点函数像什么}

可以把 \texttt{simulateTimeVaryingFirDouble} 理解成：

\begin{quote}
拿一套理想精度的计算器，完全不考虑位宽限制，去计算时变信道卷积的正确答案。
\end{quote}

\subsection{定点函数像什么}

可以把 \texttt{simulateTimeVaryingFirFixed} 理解成：

\begin{quote}
按 FPGA 的输入格式和系数量化格式做乘加，但先别急着截成 16 bit，而是把所有结果临时记在一个更宽的寄存器里，最后再统一考虑怎么截位最合理。
\end{quote}

\section{结论}

这两个函数分工明确：

\begin{itemize}
\item \texttt{simulateTimeVaryingFirDouble} 负责给出理想参考输出
\item \texttt{simulateTimeVaryingFirFixed} 负责给出定点链路下、截位前的高精度结果
\end{itemize}

后续的动态 shift 搜索、误差分析、频域比较，都是建立在这两个函数输出的基础上进行的。

因此，这两个函数可以视为整个验证工具最核心的“参考链路”和“定点链路”计算内核。

\end{document}
